\documentclass[11pt,a4paper]{article}
\usepackage[margin=1in]{geometry}
\usepackage{amsmath,amssymb,amsthm}
\usepackage{algorithm}
\usepackage{algpseudocode}
\usepackage{booktabs}
\usepackage{hyperref}

\theoremstyle{definition}
\newtheorem{definition}{Definition}[section]
\newtheorem{theorem}{Theorem}[section]
\newtheorem{lemma}[theorem]{Lemma}
\newtheorem{remark}{Remark}[section]

\title{\textbf{Rigorous Mathematical Specification, Affine Tensor Contractions, and Variance Scaling of the Linear Layer}}
\author{Team Scaibu \\ \small Email: \texttt{jaydeep.9664920749@gmail.com} \\ \small Architecture and Core Engine Team}
\date{October 2026}

\begin{document}
\maketitle

\begin{abstract}
This document presents the complete mathematical formulation and algorithmic specification of the Linear (Affine) Layer $Y = X W^T + \mathbf{1} b^T$. As the universal computational building block of deep neural networks, the affine transformation maps input manifolds across dimensional spaces. We formalize tensor contraction mechanics, establish variance preservation bounds under standard initialization schemes, provide the complete algorithmic pseudo-code, derive computational complexities, step through a numerical calculation, and address numerical considerations.
\end{abstract}

\section{Notation and Mathematical Preliminaries}

Let $X \in \mathbb{R}^{B \times d_{\text{in}}}$ denote an input mini-batch tensor with batch size $B \in \mathbb{N}_{\ge 1}$ and input dimension $d_{\text{in}} \in \mathbb{N}_{\ge 1}$.

\begin{itemize}
    \item $W \in \mathbb{R}^{d_{\text{out}} \times d_{\text{in}}}$: Trainable weight matrix mapping $\mathbb{R}^{d_{\text{in}}} \to \mathbb{R}^{d_{\text{out}}}$.
    \item $b \in \mathbb{R}^{d_{\text{out}}}$: Trainable scalar bias vector.
    \item $Y \in \mathbb{R}^{B \times d_{\text{out}}}$: Output activation batch tensor.
    \item $B, d_{\text{in}}, d_{\text{out}}$: Batch size, input feature dimension, and output feature dimension.
\end{itemize}

\begin{table}[h]
\centering
\caption{Summary of Mathematical Symbols for Linear Layer}
\begin{tabular}{lll}
\toprule
\textbf{Symbol} & \textbf{Type / Domain} & \textbf{Description} \\
\midrule
$X$ & $\mathbb{R}^{B \times d_{\text{in}}}$ & Input batch matrix \\
$W$ & $\mathbb{R}^{d_{\text{out}} \times d_{\text{in}}}$ & Linear projection weights \\
$b$ & $\mathbb{R}^{d_{\text{out}}}$ & Bias vector \\
$Y$ & $\mathbb{R}^{B \times d_{\text{out}}}$ & Output batch activations \\
\bottomrule
\end{tabular}
\end{table}

\section{Problem Definition and Core Concepts}

\subsection{Intuition and Motivation}
A linear layer performs an affine transformation consisting of rotation, scaling, reflection, and translation. In deep architectures, stacking linear transformations interleaved with element-wise non-linear activations forms a universal function approximator. For a mini-batch of $B$ independent samples, the transformation processes each sample concurrently via matrix multiplication.

\subsection{Formal Mathematical Definition}
\begin{definition}[Affine Batch Transformation]
Given input batch $X \in \mathbb{R}^{B \times d_{\text{in}}}$, weight matrix $W \in \mathbb{R}^{d_{\text{out}} \times d_{\text{in}}}$, and bias $b \in \mathbb{R}^{d_{\text{out}}}$:
\begin{equation}
Y = X W^T + \mathbf{1}_B b^T \label{eq:linear_batch}
\end{equation}
Coordinate-wise for sample $b \in \{1, \dots, B\}$ and output channel $j \in \{1, \dots, d_{\text{out}}\}$:
\begin{equation}
Y_{b, j} = \sum_{k=1}^{d_{\text{in}}} X_{b, k} W_{j, k} + b_j \label{eq:linear_elem}
\end{equation}
\end{definition}

\section{Algorithmic Specification}

The computational implementation from \texttt{NnAlgoLinearAffineLayer} is shown below.

\begin{algorithm}[H]
\caption{Linear Affine Layer Forward Evaluation}
\begin{algorithmic}[1]
\Require Input batch $X \in \mathbb{R}^{B \times d_{\text{in}}}$, weight matrix $W \in \mathbb{R}^{d_{\text{out}} \times d_{\text{in}}}$, optional bias $b \in \mathbb{R}^{d_{\text{out}}}$
\Ensure Output batch $Y \in \mathbb{R}^{B \times d_{\text{out}}}$
\State $B \gets \text{rows}(X), \quad d_{\text{in}} \gets \text{cols}(X)$
\State $d_{\text{out}} \gets \text{rows}(W)$
\If{$B = 0$ \textbf{or} $\text{cols}(W) \ne d_{\text{in}}$}
    \State \textbf{throw} PreconditionFailureException(``Matrix dimension mismatch'')
\EndIf
\State $Y \gets \text{new Matrix of shape } (B, d_{\text{out}})$
\For{$i \gets 1 \textbf{ to } B$}
    \For{$j \gets 1 \textbf{ to } d_{\text{out}}$}
        \State $\text{sum} \gets 0.0$
        \For{$k \gets 1 \textbf{ to } d_{\text{in}}$}
            \State $\text{sum} \gets \text{sum} + X[i][k] \cdot W[j][k]$
        \EndFor
        \If{$b \text{ is not null}$}
            \State $\text{sum} \gets \text{sum} + b[j]$
        \EndIf
        \State $Y[i][j] \gets \text{sum}$
    \EndFor
\EndFor
\State \Return $\{\text{output\_batch}: Y\}$
\end{algorithmic}
\end{algorithm}

\section{Theoretical Guarantees and Variance Preservation}

\begin{theorem}[Forward Output Variance Scaling]
Let input coordinates $X_{b, k}$ be zero-mean with variance $\sigma_x^2$, and weights $W_{j, k}$ be independent zero-mean with variance $\sigma_w^2$. If bias $b = \mathbf{0}$, then:
\begin{equation}
\mathbb{E}[Y_{b, j}] = 0 \quad \text{and} \quad \operatorname{Var}(Y_{b, j}) = d_{\text{in}} \sigma_x^2 \sigma_w^2
\end{equation}
\end{theorem}

\begin{proof}
By linearity of expectation: $\mathbb{E}[Y_{b, j}] = \sum_{k=1}^{d_{\text{in}}} \mathbb{E}[X_{b, k}] \mathbb{E}[W_{j, k}] = 0$.
Since terms in the sum are mutually independent:
\begin{equation}
\operatorname{Var}(Y_{b, j}) = \sum_{k=1}^{d_{\text{in}}} \operatorname{Var}(X_{b, k} W_{j, k}) = \sum_{k=1}^{d_{\text{in}}} \mathbb{E}[X_{b, k}^2] \mathbb{E}[W_{j, k}^2] = d_{\text{in}} \sigma_x^2 \sigma_w^2
\end{equation}
\end{proof}

\subsection{Limitations}
\begin{itemize}
    \item \textbf{Linear Expressivity}: A sequence of purely linear layers $Y = X W_1^T W_2^T \dots W_L^T$ collapses algebraically to a single linear map $W_{\text{eff}} = W_L \dots W_1$. Non-linear activation functions must follow linear layers to build deep representations.
\end{itemize}

\section{Complexity Analysis}

\subsection{Time Complexity}
\begin{itemize}
    \item \textbf{Matrix Contraction}: Computing $X W^T$ requires $B \times d_{\text{out}} \times d_{\text{in}}$ multiply-accumulate operations ($\mathcal{O}(B \cdot d_{\text{in}} \cdot d_{\text{out}})$ FLOPs).
    \item \textbf{Bias Addition}: Requires $B \cdot d_{\text{out}}$ scalar additions.
    \item \textbf{Total Time}: $\Theta(B \cdot d_{\text{in}} \cdot d_{\text{out}})$ FLOPs.
\end{itemize}

\subsection{Space Complexity}
\begin{itemize}
    \item \textbf{Output Tensor Memory}: Stores output matrix $Y$ of size $B \times d_{\text{out}}$ ($\Theta(B \cdot d_{\text{out}})$ auxiliary memory).
\end{itemize}

\section{Worked Numerical Example}

Let $B=1, d_{\text{in}}=2, d_{\text{out}}=2$:
\begin{equation}
X = \begin{bmatrix} 2.0 & 3.0 \end{bmatrix}, \quad W = \begin{bmatrix} 0.5 & -0.5 \\ 1.0 & 2.0 \end{bmatrix}, \quad b = \begin{bmatrix} 0.1 \\ -0.2 \end{bmatrix}
\end{equation}

\begin{align}
Y_{1, 1} &= 2.0(0.5) + 3.0(-0.5) + 0.1 = 1.0 - 1.5 + 0.1 = -0.40 \\
Y_{1, 2} &= 2.0(1.0) + 3.0(2.0) - 0.2 = 2.0 + 6.0 - 0.2 = 7.80
\end{align}
Output: $Y = \begin{bmatrix} -0.40 & 7.80 \end{bmatrix}$.

\section{Numerical Considerations and Edge Cases}

\begin{itemize}
    \item \textbf{BLAS / GEMM Acceleration}: In modern deep learning libraries, linear layer forward evaluations are dispatched directly to hardware-optimized BLAS Level-3 GEMM routines (\texttt{cublasSgemm} / \texttt{rocBLAS}).
    \item \textbf{Dimension Alignment}: Padded memory allocations aligned to 64-byte or 128-byte boundaries prevent cache line split penalties.
\end{itemize}

\begin{thebibliography}{9}
\bibitem{goodfellow2016} I.~Goodfellow, Y.~Bengio, and A.~Courville, \emph{Deep Learning}, MIT Press, 2016.
\bibitem{glorot2010lin} X.~Glorot and Y.~Bengio, ``Understanding the difficulty of training deep feedforward neural networks,'' in \emph{AISTATS}, 2010.
\end{thebibliography}

\end{document}
